%! AP Statistics — Unit 5, Lesson 5.1 — Group Activity (Answer Key)
\documentclass[10pt]{article}
\usepackage{apstats-article}
\usepackage{apstats-key}   % pulls in -boxes; do NOT also load -boxes
\newcommand{\TallMath}[1]{$\displaystyle #1\rule[-1.4em]{0pt}{3.2em}$}

\begin{document}

\pageheader{Unit 5, Lesson 5.1}{Group Activity --- Answer Key}
\namepartnerperiod

\vspace{0.15in}

\begin{tcolorbox}[colback=redbg, colframe=black!40, boxrule=0.6pt, arc=2mm,
    left=3mm, right=3mm, top=2mm, bottom=2mm, title=\textbf{Tier R --- Remediate}]
\textbf{Scenario A.}
\begin{enumerate}
  \item Is either mean the true mean? \ans{No} \quad ($\square$ Yes \quad \textcolor{keyred}{\textbf{$\checkmark$}} No)
  \item The values 11.2 and 10.8 are called \ans{statistics} because they come from a sample.
  \item The true mean is called a \ans{parameter}.
  \item ``The two sample means differ because of'' \ansline{sampling variability --- different random samples from the same orchard will contain different apples, producing different sample means by chance.}
\end{enumerate}
\end{tcolorbox}

\vspace{0.1in}

\begin{tcolorbox}[colback=redbg, colframe=black!40, boxrule=0.6pt, arc=2mm,
    left=3mm, right=3mm, top=2mm, bottom=2mm, title=\textbf{Tier A --- Approaching Proficiency}]
\textbf{Scenario B.}
\begin{enumerate}
  \item \ansline{Statistic: the sample mean commute time (28.4 min or 31.1 min). Parameter: the true mean commute time of all households in each neighborhood.}
  \item \ansline{The 2.7-minute difference could easily be due to sampling variability with samples of only 50 households. To decide, we would need to know more about the spread (standard deviation) of commute times and the true population sizes. A formal significance test (Unit 7) would be needed to rule out chance.}
  \item \textit{(Dot-plot should be roughly symmetric, centered near 28.4 min, with some spread.)}
  \item \ans{sampling distribution (of the sample mean)}
\end{enumerate}
\end{tcolorbox}

\vspace{0.1in}

\begin{tcolorbox}[colback=redbg, colframe=black!40, boxrule=0.6pt, arc=2mm,
    left=3mm, right=3mm, top=2mm, bottom=2mm, title=\textbf{Tier E --- Extension}]
\textbf{Scenario C.}
\begin{enumerate}
  \item \ansline{Parameter: the population mean daily sodium intake, $\mu = 3{,}400$ mg.}
  \item \ansline{Shape: approximately symmetric (bell-shaped), by the Central Limit Theorem (preview). Center: near 3,400 mg, because samples from this population should average around the population mean on average. Spread: smaller than the population SD of 800 mg, because averaging 25 values reduces variability.}
  \item \ansline{The spread (standard deviation) of the sampling distribution would decrease --- more values in each sample smooth out the extremes, so sample means cluster more tightly around the population mean.}
  \item \ansline{Sample question: ``Is the mean daily sodium intake of adults in a given region different from the national mean of 3,400 mg?'' (Other valid questions accepted.)}
\end{enumerate}
\end{tcolorbox}

\begin{teachernote}
Tier E prediction about shape may anticipate the CLT (Lesson 5.3). That's fine ---
encourage the reasoning without formally introducing the theorem yet.
\end{teachernote}

\end{document}
